\documentclass[12pt]{article}

% Fonts and spacing
\usepackage[T1]{fontenc}
\usepackage{newtxtext,newtxmath}
\usepackage{setspace}
\onehalfspacing

% Page layout
\usepackage[letterpaper,top=1.25in,bottom=1in,left=1in,right=1in]{geometry}

% Graphics & color
\usepackage{graphicx}
\usepackage{caption}
\usepackage{subcaption}
\usepackage{float}
\usepackage{booktabs}
\usepackage{siunitx}

% Code listings
\usepackage{listings}
\lstdefinestyle{mystyle}{
  language=Matlab,
  basicstyle=\ttfamily\footnotesize,
  numbers=left,
  numberstyle=\tiny,
  stepnumber=1,
  numbersep=5pt,
  frame=single,
  breaklines=true,
  showstringspaces=false,
  tabsize=4
}
\lstset{style=mystyle}

% Header
\usepackage{fancyhdr}
\pagestyle{fancy}
\fancyhf{}
\cfoot{\thepage}

% Hyperlinks (should be loaded last)
\usepackage[colorlinks=true,linkcolor=black,anchorcolor=black,citecolor=black,urlcolor=black]{hyperref}

\begin{document}

% ------------------------------ COVER PAGE -------------------------------
\begin{titlepage}
  \centering
  {\Large ECSE 512 -- Digital Signal Processing \par}
  \vspace{1em}
  {\Large McGill University, Fall 2025\par}
  \vspace{2em}
  {\huge\bfseries Adaptive Equalization for 4--QAM over a Dispersive Channel\par}
  \vspace{2em}
  {\large Term Project Report\par}
  \vspace{1em}
  {\large Author: \textbf{Frank Yin}, ID: \textbf{26079789}\par}
  \vspace{1em}
  {\large Date: \today\par}
  \vfill
  \begin{abstract}
  We implement and evaluate a complex LMS--based decision--directed adaptive FIR equalizer for 4--QAM transmission over frequency--selective channels that induce inter--symbol interference (ISI). The complete MATLAB simulation chain includes a channel model with unit energy, complex AWGN, a training sequence followed by a decision--directed phase, and performance assessment via learning curves and symbol error rate (SER). We study the impact of SNR, equalizer length, and step size. Results confirm that the LMS equalizer effectively mitigates ISI and tracks slow variations when properly tuned.
  \end{abstract}
  \vfill
\end{titlepage}
\setcounter{page}{1}

% ------------------------------ TOC --------------------------------------
\tableofcontents
\newpage

% ------------------------------ 1. INTRO ---------------------------------
\section{Introduction}
Modern digital communication receivers must cope with inter--symbol interference (ISI), which arises when the channel is time--dispersive or frequency--selective. In baseband terms, the received sequence can be written as the convolution of the transmit symbols with the channel impulse response plus additive noise. ISI distorts the constellation and increases the symbol error rate (SER) if left uncompensated. Equalization is therefore required, especially in wireline channels with multipath reflections and in wireless channels with delay spread.

Among the many equalization strategies---including zero--forcing (ZF), minimum mean--square error (MMSE), decision--feedback equalization (DFE), maximum likelihood sequence estimation (MLSE)~\cite{proakis}, and blind methods such as CMA~\cite{qureshi}---adaptive linear equalizers trained with the least mean squares (LMS) algorithm~\cite{haykin} strike an attractive balance between complexity, robustness, and performance. LMS only requires local gradient information and a step--size parameter, which makes it easy to implement in real time. This project implements a complex LMS adaptive FIR equalizer in a 4--QAM link and assesses its performance using Monte Carlo simulations.

\subsection*{Project goals}
The goals are to: (i) design and code a full simulation chain in MATLAB, (ii) implement training followed by decision--directed adaptation, (iii) evaluate performance as a function of SNR, equalizer length, and LMS step size, and (iv) document the methodology and findings in a technical report.

\subsection*{Report organization}
Section~\ref{sec:model} describes the system model and problem formulation. Section~\ref{sec:lms} reviews LMS equalization and design trade--offs. Section~\ref{sec:results} presents results and discussions. Section~\ref{sec:conclusion} concludes. A complete listing of our MATLAB code is provided in the Appendix.

% ------------------------------ 2. MODEL ---------------------------------
\section{System Model and Problem Formulation}\label{sec:model}
This section provides a comprehensive mathematical description of the adaptive equalization system, with each component corresponding to the block diagram in Fig.~\ref{fig:block} and the MATLAB implementation.

\subsection*{Transmitter and Signal Generation}
The data generator produces an independent and identically distributed (IID) sequence of information symbols $x[n]$ drawn from a 4--QAM constellation $S$. The constellation consists of four equi--spaced points in the complex plane:
\begin{equation}
S = \left\{s_1, s_2, s_3, s_4\right\} = \frac{\sigma_s}{\sqrt{2}}\left\{1+j, -1+j, -1-j, 1-j\right\},
\end{equation}
where $\sigma_s^2$ is the symbol power, ensuring $E\{|x[n]|^2\} = \sigma_s^2$. Each symbol is selected with equal probability $P(x[n]=s_i) = 1/4$ for $i=1,\ldots,4$. In our implementation, symbols are generated by mapping random indices $\{1,2,3,4\}$ to constellation points via the function \texttt{qam4\_mod()}, which directly implements the mapping $x[n] = S[\text{idx}[n]]$.

\subsection*{Channel Model and Inter--Symbol Interference}
The channel filter $h[n]$ models the combined effect of transmit filtering, physical propagation (multipath, reflections), and receive filtering. It is represented as a finite impulse response (FIR) filter of length $K$:
\begin{equation}
h[n] = \begin{cases} h_n & \text{for } n=0,1,\ldots,K-1 \\ 0 & \text{otherwise} \end{cases}
\end{equation}
The channel is normalized to unit energy such that $\sum_{n=0}^{K-1}|h[n]|^2 = 1$, ensuring that the channel does not amplify or attenuate the average signal power. This normalization is implemented via \texttt{normalize\_channel()}, which computes $h \leftarrow h/\|h\|_2$.

Our simulation supports three channel types, implemented in \texttt{generate\_channel()}:
\begin{itemize}
\item \textbf{Mild ISI:} A 3--tap channel $h = [0.9+0j, 0.3-0.2j, 0.15+0.05j]$ with modest postcursor interference, representing a relatively benign dispersive channel.
\item \textbf{Severe ISI:} A 5--tap channel $h = [0.3+0.2j, 0.9+0j, 0.25-0.15j, 0.15+0.10j, 0.05-0.06j]$ with both precursor and postcursor interference, modeling strong multipath effects.
\item \textbf{Random Rayleigh:} A $K$--tap channel with complex Gaussian taps $h[k] \sim \mathcal{CN}(0, \sigma_h^2)$ before normalization, representing a frequency--selective fading channel.
\end{itemize}

The channel output before noise addition is the convolution:
\begin{equation}
y_{\text{clean}}[n] = \sum_{k=0}^{K-1} h[k]\,x[n-k] = (h*x)[n],
\end{equation}
which introduces inter--symbol interference (ISI) when $K > 1$, as each received sample contains contributions from multiple transmitted symbols.

\subsection*{Additive Noise Model}
The noise generator produces complex circularly symmetric additive white Gaussian noise (AWGN) $w[n] \sim \mathcal{CN}(0, \sigma_n^2)$, where $\sigma_n^2$ is chosen to achieve a target signal--to--noise ratio (SNR) defined as:
\begin{equation}
\text{SNR}_{\text{dB}} = 10\log_{10}\left(\frac{\sigma_s^2}{\sigma_n^2}\right).
\end{equation}
Since the channel is unit energy, $E\{|y_{\text{clean}}[n]|^2\} \approx \sigma_s^2$, and the noise variance is computed as $\sigma_n^2 = \sigma_s^2 / 10^{\text{SNR}_{\text{dB}}/10}$. The function \texttt{add\_awgn()} generates noise samples as $w[n] = \sqrt{\sigma_n^2/2}(w_r[n] + jw_i[n])$ where $w_r[n], w_i[n] \sim \mathcal{N}(0,1)$ are independent real Gaussian random variables.

The final received signal at the equalizer input is:
\begin{equation}
y[n] = y_{\text{clean}}[n] + w[n] = \sum_{k=0}^{K-1} h[k]\,x[n-k] + w[n].
\end{equation}
This signal exhibits both ISI distortion and noise corruption, making direct symbol detection unreliable.

\subsection*{Adaptive Equalizer Structure}
The adaptive equalizer is an $N$--tap complex FIR filter with time--varying coefficients $\mathbf{c}[n] = [c_0[n], c_1[n], \ldots, c_{N-1}[n]]^T$. At each time $n$, it forms a linear combination of the $N$ most recent received samples:
\begin{equation}
\hat{x}[n] = \sum_{k=0}^{N-1} c_k[n]\,y[n-k] = \mathbf{c}^H[n]\,\mathbf{y}[n],
\end{equation}
where $\mathbf{y}[n] = [y[n], y[n-1], \ldots, y[n-N+1]]^T$ is the input regressor vector. The equalizer output $\hat{x}[n]$ is a soft estimate of the transmitted symbol, which must be processed by a decision device to produce a hard decision.

\subsection*{Decision Device and Symbol Detection}
The decision device implements minimum--distance detection, mapping the equalizer output $\hat{x}[n]$ to the nearest constellation point:
\begin{equation}
\tilde{x}[n] = \mathcal{Q}(\hat{x}[n]) = \arg\min_{s \in S} |\hat{x}[n] - s|^2.
\end{equation}
This is implemented in \texttt{qam4\_demod()} by computing distances to all four constellation points and selecting the index $k$ that minimizes $|\hat{x}[n] - S[k]|^2$, then outputting $\tilde{x}[n] = S[k]$.

\subsection*{Training and Decision--Directed Modes}
The system operates in two distinct phases, controlled by a switch in the block diagram (Fig.~\ref{fig:block}):

\textbf{Training Phase ($n \leq L_{\text{tr}}$):} During the initial $L_{\text{tr}}$ symbols, the transmitter sends a known training sequence. The desired signal for the equalizer is set to $d[n] = x[n-D]$, where $D$ is the \emph{decision delay}. This delay accounts for the fact that the equalizer output $\hat{x}[n]$ estimates a symbol transmitted $D$ samples earlier due to the channel's memory. The delay $D$ is chosen to align the equalizer's center tap with the dominant channel tap, as computed by \texttt{pick\_decision\_delay()}, which selects $D = \arg\max_k |h[k]| - 1$ (clipped to $[0, N-1]$).

\textbf{Decision--Directed Phase ($n > L_{\text{tr}}$):} After training, the system switches to decision--directed mode, where the desired signal becomes $d[n] = \tilde{x}[n] = \mathcal{Q}(\hat{x}[n])$. This allows the equalizer to continue adapting to track slow channel variations without requiring additional training overhead. The switch occurs at index $n = N-1 + L_{\text{tr}}$ in our implementation, ensuring that $N$ received samples are available before adaptation begins.

\subsection*{Error Signal and Adaptation}
The error signal is computed as the difference between the desired and actual equalizer output:
\begin{equation}
e[n] = d[n] - \hat{x}[n],
\end{equation}
where $d[n]$ switches between $x[n-D]$ (training) and $\tilde{x}[n]$ (decision--directed) as described above. This error signal drives the LMS adaptation algorithm, which updates the equalizer coefficients to minimize the mean squared error $E\{|e[n]|^2\}$.

\subsection*{Performance Metrics}
We evaluate system performance using two primary metrics:

\textbf{Learning Curve:} The instantaneous squared error $|e[n]|^2$ versus time $n$, which reveals the convergence behavior during training and the steady--state residual error. The learning curve is computed directly in \texttt{lms\_equalizer()} and averaged over multiple Monte Carlo trials to estimate $E\{|e[n]|^2\}$.

\textbf{Symbol Error Rate (SER):} The probability of symbol detection error, computed as:
\begin{equation}
\text{SER} = \frac{1}{L_{\text{data}}} \sum_{n=n_0}^{L} \mathbf{1}\{\tilde{x}[n] \neq x[n-D]\},
\end{equation}
where $n_0$ marks the start of the decision--directed phase and $L_{\text{data}}$ is the number of data symbols. The SER is measured only during decision--directed operation (implemented in \texttt{measure\_ser()}) and averaged over multiple independent realizations to obtain reliable statistics.

\begin{figure}[H]
 \centering
 \includegraphics[width=\textwidth]{figures/block_diagram.png}
 \caption{Block diagram of the data transmission system with adaptive equalization. The system processes transmitted symbols $x[n]$ through a dispersive channel $h[n]$ with additive noise $w[n]$, producing the received signal $y[n]$. The adaptive equalizer with coefficients $c_k[n]$ generates soft estimates $\hat{x}[n]$, which are quantized by the decision device to produce hard decisions $\tilde{x}[n]$. The error signal $e[n] = d[n] - \hat{x}[n]$ drives the LMS adaptation, where $d[n]$ switches between the delayed training sequence $x[n-D]$ and the decision--directed reference $\tilde{x}[n]$.}
 \label{fig:block}
\end{figure}

% ------------------------------ 3. LMS -----------------------------------
\section{Adaptive LMS Equalization}\label{sec:lms}
This section provides a detailed review of the least mean squares (LMS) algorithm for adaptive equalization, including its mathematical derivation, implementation details, and design trade--offs for parameter selection.

\subsection*{LMS Algorithm Derivation}
The LMS algorithm is a stochastic gradient descent method that minimizes the mean squared error (MSE) cost function:
\begin{equation}
J(\mathbf{c}) = E\{|e[n]|^2\} = E\{|d[n] - \mathbf{c}^H\mathbf{y}[n]|^2\},
\end{equation}
where $\mathbf{y}[n] = [y[n], y[n-1], \ldots, y[n-N+1]]^T$ is the input regressor vector and $\mathbf{c} = [c_0, c_1, \ldots, c_{N-1}]^T$ is the equalizer coefficient vector. The optimal Wiener solution is $\mathbf{c}_{\text{opt}} = \mathbf{R}_y^{-1}\mathbf{p}_{yd}$, where $\mathbf{R}_y = E\{\mathbf{y}[n]\mathbf{y}^H[n]\}$ is the input correlation matrix and $\mathbf{p}_{yd} = E\{\mathbf{y}[n]d^*[n]\}$ is the cross--correlation vector.

Since computing $\mathbf{R}_y^{-1}$ requires knowledge of channel statistics, LMS uses an instantaneous gradient estimate. The true gradient of $J(\mathbf{c})$ with respect to $\mathbf{c}^*$ is:
\begin{equation}
\nabla_{\mathbf{c}^*} J = -E\{\mathbf{y}[n]e^*[n]\},
\end{equation}
where $e[n] = d[n] - \mathbf{c}^H[n]\mathbf{y}[n]$. LMS approximates this by the instantaneous gradient:
\begin{equation}
\hat{\nabla}_{\mathbf{c}^*} J[n] = -\mathbf{y}[n]e^*[n],
\end{equation}
leading to the update equation:
\begin{equation}
\mathbf{c}[n+1] = \mathbf{c}[n] - \mu \hat{\nabla}_{\mathbf{c}^*} J[n] = \mathbf{c}[n] + \mu \mathbf{y}[n]e^*[n],
\end{equation}
where $\mu > 0$ is the step size parameter. This is the complex LMS algorithm~\cite{haykin}, which adapts the equalizer coefficients at each time step using only the current input vector and error signal.

\subsection*{Algorithm Implementation}
The complete LMS equalization algorithm, as implemented in \texttt{lms\_equalizer()}, proceeds as follows:

\textbf{Initialization:} The equalizer coefficients are initialized as a delayed delta function:
\begin{equation}
c_k[0] = \begin{cases} 1 & \text{if } k = D \\ 0 & \text{otherwise} \end{cases},
\end{equation}
where $D$ is the decision delay. This initialization, implemented as \texttt{c(min(D+1,N)) = 1}, places the equalizer's main tap at position $D$, providing a reasonable starting point that forwards the $D$--th delayed sample. This accelerates initial convergence compared to zero initialization.

\textbf{Main Loop:} For each time index $n \geq N$ (ensuring all required input samples are available):
\begin{enumerate}
\item \textbf{Form input vector:} $\mathbf{y}[n] = [y[n], y[n-1], \ldots, y[n-N+1]]^T$, implemented as \texttt{flipud(y(n-N+1:n))}.
\item \textbf{Compute equalizer output:} $\hat{x}[n] = \mathbf{c}^H[n]\mathbf{y}[n]$, using the current coefficient vector.
\item \textbf{Determine desired signal:}
\begin{equation}
d[n] = \begin{cases}
x[n-D] & \text{if } n \leq L_{\text{tr}} \text{ (training mode)} \\
\mathcal{Q}(\hat{x}[n]) & \text{if } n > L_{\text{tr}} \text{ (decision--directed mode)}
\end{cases}
\end{equation}
where $\mathcal{Q}(\cdot)$ denotes the minimum--distance quantizer to the 4--QAM constellation.
\item \textbf{Compute error:} $e[n] = d[n] - \hat{x}[n]$.
\item \textbf{Update coefficients:} $\mathbf{c}[n+1] = \mathbf{c}[n] + \mu \mathbf{y}[n]e^*[n]$.
\item \textbf{Store metrics:} Record $|e[n]|^2$ for the learning curve and $\mathbf{c}[n]$ for coefficient evolution analysis.
\end{enumerate}

The algorithm maintains a history of coefficients $\mathbf{c}_{\text{hist}}[n]$ and error signals, enabling post--processing analysis of convergence behavior.

\subsection*{Convergence Analysis}
The convergence properties of LMS depend critically on the step size $\mu$. For convergence in the mean, the step size must satisfy:
\begin{equation}
0 < \mu < \frac{2}{\lambda_{\max}},
\end{equation}
where $\lambda_{\max}$ is the largest eigenvalue of $\mathbf{R}_y$. A more practical bound uses the trace of $\mathbf{R}_y$:
\begin{equation}
0 < \mu < \frac{2}{\mathrm{tr}\{\mathbf{R}_y\}} = \frac{2}{N P_y},
\end{equation}
where $P_y = E\{|y[n]|^2\}$ is the average received power. For a unit--energy channel, $P_y \approx \sigma_s^2 + \sigma_n^2$, leading to the rule of thumb:
\begin{equation}
\mu \lesssim \frac{1}{N(\sigma_s^2 + \sigma_n^2)}.
\end{equation}

The convergence rate is approximately exponential with time constant $\tau \approx 1/(\mu \lambda_{\min})$ for small $\mu$, where $\lambda_{\min}$ is the smallest eigenvalue of $\mathbf{R}_y$. The steady--state excess MSE (misadjustment) is:
\begin{equation}
M = \frac{\mu}{2} \mathrm{tr}\{\mathbf{R}_y\} = \frac{\mu N P_y}{2},
\end{equation}
which increases linearly with $\mu$ and $N$. This creates a fundamental trade--off: larger $\mu$ accelerates convergence but increases steady--state error, while smaller $\mu$ reduces misadjustment but slows convergence.

\subsection*{Parameter Selection and Design Trade--Offs}

\textbf{Equalizer Length $N$:} The equalizer must span the channel's delay spread to effectively cancel ISI. If the channel has significant energy over $K$ taps, the equalizer length should satisfy $N \gtrsim K$ to capture the channel memory. However, choosing $N$ too large introduces several drawbacks:
\begin{itemize}
\item \textbf{Increased complexity:} Each additional tap requires one multiply--accumulate operation per iteration.
\item \textbf{Slower convergence:} More parameters to adapt increases the convergence time, approximately proportional to $N$.
\item \textbf{Higher misadjustment:} The excess MSE scales as $M \propto \mu N P_y$, so longer equalizers accumulate more steady--state error for fixed $\mu$.
\item \textbf{Over--fitting:} Excessively long equalizers may adapt to noise rather than signal structure.
\end{itemize}
In practice, $N$ should be chosen as the smallest value that spans the significant channel energy, typically $N = K + 2$ to $N = 2K$ to provide margin for adaptation. Our simulations use $N=11$ for a 5--tap severe channel, providing adequate span without excessive complexity.

\textbf{Step Size $\mu$:} The step size is the most critical parameter, balancing convergence speed and steady--state performance. The selection process involves several considerations:

\begin{itemize}
\item \textbf{Convergence speed:} Larger $\mu$ reduces the convergence time constant, allowing the equalizer to adapt quickly during the training phase. This is particularly important when $L_{\text{tr}}$ is limited.
\item \textbf{Steady--state error:} Smaller $\mu$ reduces misadjustment, leading to lower residual MSE and better SER performance after convergence.
\item \textbf{Stability:} Very large $\mu$ can cause instability, especially in decision--directed mode where error propagation can amplify coefficient fluctuations.
\item \textbf{Training duration:} If $L_{\text{tr}}$ is short, a larger $\mu$ may be necessary to achieve convergence within the training window. Conversely, if training is long, smaller $\mu$ can be used to achieve lower steady--state error.
\end{itemize}

Our implementation uses $\mu = 0.003$ as a default, which provides a good balance for $N=11$ and $L_{\text{tr}}=1500$. This corresponds to $\mu \approx 0.27/(N P_y)$ for typical SNR values, well within the stability bound.

\textbf{Decision Delay $D$:} The decision delay aligns the desired symbol $x[n-D]$ with the equalizer's estimate $\hat{x}[n]$. The optimal $D$ maximizes the signal--to--interference ratio at the equalizer output. A heuristic approach, implemented in \texttt{pick\_decision\_delay()}, selects $D$ to align with the dominant channel tap:
\begin{equation}
D = \arg\max_{k} |h[k]| - 1,
\end{equation}
clipped to the range $[0, N-1]$. This ensures that the equalizer's center tap (at index $D$) targets the strongest channel path, facilitating effective ISI cancellation. For the severe 5--tap channel with dominant tap at index 1 (0--indexed), this yields $D=1$, placing the main equalizer tap at the second position.

\textbf{Training Length $L_{\text{tr}}$:} The training duration must be sufficient for the equalizer to converge. The required length depends on the convergence time constant $\tau \approx 1/(\mu \lambda_{\min})$ and the desired convergence factor. A rule of thumb is $L_{\text{tr}} \gtrsim 5\tau$ to achieve convergence within 1\% of steady--state. For our parameters ($\mu=0.003$, $N=11$), this suggests $L_{\text{tr}} \gtrsim 1000$ symbols, and we use $L_{\text{tr}} = 1500$ to provide margin. Longer training improves initial convergence but reduces the fraction of time available for data transmission.

\textbf{Training vs Decision--Directed Mode:} The two--phase operation addresses different challenges:
\begin{itemize}
\item \textbf{Training phase:} Uses the known transmitted sequence $x[n-D]$ as the desired signal, avoiding error propagation and ensuring reliable convergence. This is essential during initial acquisition when the equalizer coefficients are far from optimal.
\item \textbf{Decision--directed phase:} Uses the quantized equalizer output $\tilde{x}[n]$ as the desired signal, enabling continuous adaptation without training overhead. This is effective once the equalizer has converged and the decision error rate is low, allowing tracking of slow channel variations.
\end{itemize}
The transition occurs automatically at $n = N-1 + L_{\text{tr}}$ in our implementation, ensuring seamless switching once sufficient training has occurred.

\subsection*{Implementation Considerations}
Several practical details in our MATLAB implementation enhance algorithm robustness:

\textbf{Boundary handling:} The algorithm begins adaptation at $n = N$ to ensure all required input samples $y[n-k]$ for $k=0,\ldots,N-1$ are available. The first $N-1$ outputs are set to zero.

\textbf{Delay alignment:} When accessing $x[n-D]$ during training, bounds checking ensures $n-D \in [1, L]$ to avoid array out--of--bounds errors. If the delay places the reference outside the transmitted sequence, $d[n]$ is set to zero.

\textbf{Constellation scaling:} The decision device uses the same constellation scaling $\sigma_s$ as the modulator, ensuring consistent symbol mapping. The function \texttt{qam4\_demod()} computes distances using the normalized constellation points.

\textbf{Error tracking:} The algorithm maintains a complete history of coefficients $\mathbf{c}_{\text{hist}}[n]$ and error signals $e[n]$, enabling detailed analysis of convergence dynamics and coefficient evolution over time.

These implementation details ensure numerical stability and accurate simulation of the adaptive equalization system, providing reliable performance metrics for system evaluation.

% ------------------------------ 4. RESULTS --------------------------------
\section{Results and Discussion}\label{sec:results}
All results were produced with the provided MATLAB code. Unless stated otherwise, the default configuration uses an $N=11$ tap equalizer, step size $\mu=0.003$, a ``severe'' 5--tap channel normalized to unit energy, a training length of $1500$ symbols followed by $15000$ payload symbols, and decision delay chosen automatically from the dominant channel tap.

\subsection*{Learning curve}
Figure~\ref{fig:learning} shows the learning curve for a single run at \SI{20}{dB} SNR, displaying the mean squared error (MSE) $|e[n]|^2$ in decibels as a function of time index $n$. The curve is smoothed using a 50--sample moving average to reduce fluctuations. The plot reveals several key characteristics: (i) rapid initial convergence during the training phase, with the MSE dropping by approximately 15--20 dB within the first 500 symbols, (ii) a transition point marked by the vertical dashed line at $n=1500$ where the system switches from training to decision--directed mode, and (iii) a steady--state residual error level (horizontal dashed line) computed as the average of the last 100 samples. 

The convergence rate is consistent with the theoretical time constant $\tau \approx 1/(\mu \lambda_{\min})$, where $\lambda_{\min}$ is the smallest eigenvalue of the input correlation matrix. The steady--state MSE reflects the combined effects of channel noise ($\sigma_n^2$), residual ISI from imperfect equalization, and LMS misadjustment $M = \mu N P_y/2 \approx 0.05$ (in linear scale). The smooth transition at the training/decision--directed boundary indicates that the equalizer has converged sufficiently to maintain stable operation in decision--directed mode, validating the choice of $L_{\text{tr}} = 1500$ symbols.

\begin{figure}[H]
 \centering
 \includegraphics[width=0.85\textwidth]{../results/learning_curve_SNR20dB.png}
 \caption{Learning curve showing MSE convergence for a single run at \SI{20}{dB} SNR. The vertical dashed line indicates the transition from training to decision--directed mode at $n=1500$. The horizontal dashed line shows the steady--state MSE level.}
 \label{fig:learning}
\end{figure}

Figure~\ref{fig:constellations} presents constellation diagrams before and after equalization. The left subplot shows the received signal $y[n]$ at the equalizer input, where the ISI and noise have spread the four ideal 4--QAM points into overlapping clusters with significant dispersion. The right subplot displays the equalizer output $\hat{x}[n]$, where the symbols have been successfully compressed back toward the ideal constellation points with minimal residual spread. The reduction in cluster size directly corresponds to the ISI cancellation achieved by the adaptive filter. The residual dispersion in the equalized constellation is primarily due to the noise floor (which cannot be removed by linear equalization) and the steady--state misadjustment. This visual confirmation demonstrates that the LMS equalizer effectively mitigates ISI and restores the signal quality necessary for reliable symbol detection, with the decision regions clearly separated.

\begin{figure}[H]
 \centering
 \includegraphics[width=\textwidth]{../results/constellations_one_run_SNR20dB.png}
 \caption{Constellation diagrams before (left) and after (right) equalization at \SI{20}{dB} SNR. The equalizer output shows successful compression of the distorted clusters back toward the ideal 4--QAM constellation points.}
 \label{fig:constellations}
\end{figure}

\subsection*{SER versus SNR}
Figure~\ref{fig:ser_snr} presents the symbol error rate (SER) as a function of SNR, obtained through Monte Carlo averaging over 5 independent trials for each SNR point. The results span SNRs from \SI{5}{dB} to \SI{30}{dB} in \SI{5}{dB} increments. As expected, the SER decreases monotonically with increasing SNR, following an approximately exponential trend on the logarithmic scale with a slope of roughly 1--2 dB per decade in the waterfall region (\SI{5}{dB}--\SI{20}{dB}). 

The curve exhibits a characteristic waterfall region at low to moderate SNRs, where noise dominates performance and the SER improves by approximately an order of magnitude for every 10 dB increase in SNR. Above \SI{25}{dB}, the curve transitions to an error floor where residual ISI and equalizer misadjustment become the limiting factors, preventing further improvement despite increasing SNR. The error floor at approximately $\text{SER} \approx 10^{-3}$ is consistent with the theoretical misadjustment $M = \mu N P_y/2$ and residual ISI from the 5--tap severe channel. For the configuration shown ($N=11$, $\mu=0.003$), the SER reaches approximately $10^{-3}$ at \SI{30}{dB} SNR, demonstrating effective ISI mitigation while highlighting the fundamental limitation of linear equalization in high--SNR regimes.

\begin{figure}[H]
 \centering
 \includegraphics[width=0.75\textwidth]{../results/SER_vs_SNR.png}
 \caption{Symbol error rate versus SNR for the adaptive LMS equalizer. Results are averaged over 5 Monte Carlo trials. The equalizer parameters are $N=11$ and $\mu=0.003$.}
 \label{fig:ser_snr}
\end{figure}

\subsection*{Impact of $\mu$ and $N$}
The step size $\mu$ plays a critical role in balancing convergence speed and steady--state performance. Figure~\ref{fig:ser_mu} shows the SER as a function of $\mu$ at a fixed SNR of \SI{20}{dB}. The results reveal a clear trade--off: for very small step sizes ($\mu < 0.001$), the equalizer converges too slowly and may not fully adapt within the 1500--symbol training period, leading to suboptimal performance with SER $\approx 10^{-2}$. As $\mu$ increases, the SER improves up to an optimal range around $\mu \approx 0.002$--$0.003$, where the equalizer achieves good convergence within the training window while maintaining low misadjustment, reaching SER $\approx 2 \times 10^{-3}$. Beyond this range, larger step sizes ($\mu > 0.005$) increase the steady--state misadjustment according to $M = \mu N P_y/2$, causing the SER to degrade back to $\approx 5 \times 10^{-3}$ at $\mu = 0.01$. The optimal $\mu$ value depends on the channel characteristics, equalizer length, and training duration, with the observed optimum aligning with the theoretical guideline $\mu \approx 1/(N P_y)$.

\begin{figure}[H]
 \centering
 \includegraphics[width=0.75\textwidth]{../results/SER_vs_mu.png}
 \caption{Symbol error rate versus LMS step size $\mu$ at \SI{20}{dB} SNR. The equalizer length is $N=11$. Results demonstrate the trade--off between convergence speed and steady--state misadjustment.}
 \label{fig:ser_mu}
\end{figure}

The equalizer length $N$ determines the span over which the equalizer can compensate for ISI. Figure~\ref{fig:ser_n} illustrates the SER as a function of $N$ at \SI{20}{dB} SNR with $\mu=0.003$. For very short equalizers ($N=7$), the SER is relatively high (approximately $8 \times 10^{-3}$) because the equalizer cannot span the full 5--tap channel memory, leaving residual ISI from the unspanned taps. As $N$ increases to $N=9$ and $N=11$, the SER improves significantly (to approximately $2 \times 10^{-3}$) as the equalizer gains sufficient span to mitigate the ISI, with $N=11$ providing approximately 2.2$\times$ the channel length. However, further increasing $N$ to $13$ and $15$ shows diminishing returns and may even slightly degrade performance due to increased misadjustment $M \propto N$ from the larger number of adaptive parameters. The optimal $N$ should be chosen to slightly exceed the channel memory length while avoiding unnecessary complexity. For the 5--tap channel used here, $N=11$ provides a good balance between ISI cancellation and misadjustment, consistent with the design guideline $N \approx 2K$.

\begin{figure}[H]
 \centering
 \includegraphics[width=0.75\textwidth]{../results/SER_vs_N.png}
 \caption{Symbol error rate versus equalizer length $N$ at \SI{20}{dB} SNR with $\mu=0.003$. The results show that $N$ should be chosen to span the channel memory while avoiding excessive complexity.}
 \label{fig:ser_n}
\end{figure}

% ------------------------------ 5. CONCLUSION -----------------------------
\section{Conclusion}\label{sec:conclusion}
We developed a complete MATLAB framework for 4--QAM transmission over an ISI channel with an adaptive LMS linear equalizer. The system successfully implements training followed by decision--directed tracking, demonstrating effective ISI mitigation across a range of channel conditions and parameter settings.

The simulation results validate the theoretical predictions: (i) the learning curve exhibits exponential convergence consistent with LMS theory, reaching steady--state within the training period, (ii) constellation diagrams confirm successful ISI cancellation with clear separation of decision regions, (iii) SER versus SNR shows the expected waterfall region transitioning to an error floor at high SNR, limited by residual ISI and misadjustment, and (iv) parameter sweeps reveal optimal operating points for $\mu$ and $N$ that balance convergence speed, steady--state error, and complexity.

Key findings include: the optimal step size $\mu \approx 0.003$ aligns with the theoretical guideline $\mu \approx 1/(N P_y)$, the equalizer length $N=11$ provides effective ISI cancellation for the 5--tap channel while avoiding excessive misadjustment, and the error floor at $\text{SER} \approx 10^{-3}$ reflects the fundamental limitations of linear equalization in high--SNR regimes. The two--phase operation (training + decision--directed) enables reliable convergence and continuous tracking without excessive overhead.

The framework provides a solid foundation for further extensions, including decision--feedback equalization (DFE) to reduce error floors, normalized LMS (NLMS) for improved convergence, MMSE training for better initial conditions, or blind equalization algorithms such as constant modulus algorithm (CMA) for scenarios without training sequences.

% ------------------------------ REFERENCES --------------------------------
\begin{thebibliography}{9}
\bibitem{haykin} S.~Haykin, \emph{Adaptive Filter Theory}, 4th~ed., Prentice Hall, 2002.
\bibitem{proakis} J.~G.~Proakis, \emph{Digital Communications}, McGraw--Hill, 2001.
\bibitem{lucky} R.~W.~Lucky, ``Techniques for adaptive equalization of digital communication,'' \emph{Bell Syst.\ Tech.\ J.}, vol.~45, pp.~255--286, Feb.~1966.
\bibitem{qureshi} S.~U.~H.~Qureshi, ``Adaptive equalization,'' \emph{Proc.\ IEEE}, vol.~73, pp.~1349--1387, Sept.~1985.
\end{thebibliography}

% ------------------------------ APPENDIX ----------------------------------
\appendix
\section*{Appendix: MATLAB Listings}
All MATLAB source files are reproduced below for completeness. Paths are relative to the \texttt{matlab/} directory.

\subsection*{main.m}\lstinputlisting{../matlab/main.m}
\subsection*{simulate\_one\_run.m}\lstinputlisting{../matlab/simulate_one_run.m}
\subsection*{lms\_equalizer.m}\lstinputlisting{../matlab/lms_equalizer.m}
\subsection*{qam4\_constellation.m}\lstinputlisting{../matlab/qam4_constellation.m}
\subsection*{qam4\_mod.m}\lstinputlisting{../matlab/qam4_mod.m}
\subsection*{qam4\_demod.m}\lstinputlisting{../matlab/qam4_demod.m}
\subsection*{generate\_channel.m}\lstinputlisting{../matlab/generate_channel.m}
\subsection*{normalize\_channel.m}\lstinputlisting{../matlab/normalize_channel.m}
\subsection*{pick\_decision\_delay.m}\lstinputlisting{../matlab/pick_decision_delay.m}
\subsection*{add\_awgn.m}\lstinputlisting{../matlab/add_awgn.m}
\subsection*{plot\_constellation.m}\lstinputlisting{../matlab/plot_constellation.m}
\subsection*{measure\_ser.m}\lstinputlisting{../matlab/measure_ser.m}

\end{document}
